\section{Results}
\label{sec:results}

Results are organized along the divergence evidence each factor in
\S\ref{sec:data} was selected to demonstrate, not alphabetically: an anchor
case, a case where the two teams agree, two cases where they reach opposite
verdicts, a case of divergence attributable to data vintage rather than
method, and (tentatively) a case where both teams agree the factor fails.
Cross-factor material -- the RQ2 autonomy results, the gap-decomposition
summary, the flagship sensitivity-band figure, and the outcome
classification -- is reported once, in a single table or figure spanning
every factor, rather than repeated inside each per-factor subsection.

\subsection{RQ2: agent autonomy and capability boundary (cross-factor)}

Table~\ref{tab:autonomy} reports, for every factor, the fraction of
high-impact MethodSpec fields the agent resolved without a human
correction, the end-to-end completion rate (whether the pipeline reached an
executable $\Aref$ with zero manual code edits), and where human correction
concentrated -- split into purely technical corrections (a type mismatch, an
out-of-menu value) versus corrections that supplied a genuine empirical
judgment the paper itself left unstated. This second split matters for
\S\ref{sec:discussion}: only the latter category threatens the leakage-proof
boundary in \S\ref{sec:framework}, since a technical correction carries no
information about how the underlying factor should be implemented.
% T1: RQ2 result -- per-factor agent autonomy footprint / capability boundary.
% Source: SourcedValue.evidence "human correction: {reason}" entries in the
% resolved MethodSpec + comparison.json's spec_quality section.
\begin{table}[htbp]
  \centering
  \caption{Agent autonomy footprint and capability boundary, by factor.}
  \label{tab:autonomy}
  \begin{tabular}{lccccc}
    \toprule
    Factor & High-impact fields & Human-corrected & Technical-only & Empirical-judgment & Zero-edit e2e? \\
    \midrule
    AssetGrowth & TODO & TODO & TODO & TODO & TODO \\
    GP          & TODO & TODO & TODO & TODO & TODO \\
    PS          & TODO & TODO & TODO & TODO & TODO \\
    BrandInvest & TODO & TODO & TODO & TODO & TODO \\
    OperProfRD  & TODO & TODO & TODO & TODO & TODO \\
    grcapx3y    & TODO & TODO & TODO & TODO & TODO \\
    \bottomrule
  \end{tabular}
\end{table}


\subsection{Per-factor deep dives}

\subsubsection{AssetGrowth (anchor case)}

AssetGrowth \citep{CooperGulenSchill2008} is the already-run anchor
case (13 tracks) and is used here to walk through how the paper-anchored
three-term identity and the field-level gap decomposition are read before
the remaining factors, each of which isolates a narrower piece of the same
picture, are discussed more briefly.

\subsubsection{GP: a rare case of agreement}

Gross profitability \citep{NovyMarx2013} is one of the few factors in the
sample where Hou-Xue-Zhang and Chen-Zimmermann's own reported numbers do
not disagree in direction or in economic magnitude. Its role in this paper
is as a control: if the paper-anchored decomposition (\S\ref{sec:framework})
attributes GP's small residual gap to configuration noise rather than to a
one-sided departure from the paper, that is itself evidence that the
decomposition does not manufacture disagreement where none exists.

\subsubsection{PS: a clean, opposite verdict}

Piotroski's \citep{Piotroski2000} F-score is the case introduced in \S\ref{sec:intro}: \citet{HouXueZhang2020} report it fails under their standardized protocol,
while Chen and Zimmermann, using the paper's own stated method, replicate it
at $t=3.29$. This factor also carries a declared, disclosed deviation: the
paper's original test restricts to the highest book-to-market quintile, a
subset condition this engine's universe-filter machinery cannot express
(\S\ref{sec:scope}), so the reported $\Aref$ is a full-sample F-score sort
rather than the paper's high-BM-subset sort, marked
\texttt{accepted\_unapplied} rather than silently approximated.

\subsubsection{fgr5yrLag: divergence this framework cannot attribute}

La Porta's \citep{LaPorta1996} long-term analyst growth-forecast factor is included
specifically because its HXZ-versus-C\&Z divergence is not, on the
available evidence, explainable by a portfolio-construction choice at all:
it is consistent instead with the two teams having pulled IBES consensus
forecast data at different vintages, a database that is itself revised
retroactively. This is stated as a limitation of the present design, not
downplayed: the paper-anchored identity in \S\ref{sec:framework} attributes
gaps to signal, configuration, or agent residual, and has no term for "the
underlying database changed between two pulls." fgr5yrLag is reported here
precisely to make that gap in the framework visible rather than to resolve
it (\S\ref{sec:discussion}).

\subsubsection{OScore: severe divergence, a known engineering gap}

Ohlson's \citep{Ohlson1980} O-score, as implemented by \citet{Dichev1998}, is a second case
where Hou, Xue, and Zhang report failure and Chen and Zimmermann report
success under the paper's own method -- more severe than PS's gap. Its
signal is a straightforward logit combination of Compustat fields, so the
divergence is unlikely to be a data-availability issue; what it does
require, and what this engine does not yet support, is the paper's
asymmetric long-low-70\%/short-high-10\% portfolio split, which cannot be
expressed by an equal-quantile breakpoint menu (\S\ref{sec:scope}). Results
reported for OScore in this section are therefore [TODO: state explicitly,
once run, whether OScore's numbers reflect the paper's asymmetric split or
an equal-decile approximation, and flag accordingly].

\subsubsection{FailureProbability: both teams report failure (tentative)}

% TODO (2026-08-19): inclusion of this factor is still undecided pending its
% engineering cost -- see docs/paper-outline.md \S5.1 (mktrf as a signal
% input, ltq/cheq quarterly fields, daily volatility aggregation). If
% dropped, remove this subsection and its rows from every cross-factor
% table/figure below and adjust the factor count throughout the paper.
Campbell, Hilscher, and Szilagyi's \citep{CampbellHilscherSzilagyi2008} failure probability is the
sample's second control case: both teams report that it fails to predict
returns under their respective protocols. Together with GP, it anchors the
sensitivity band in \S\ref{sec:sensitivity} at both ends -- a factor neither
team's convention can rescue, and a factor neither team's convention can
break.

\subsection{Cross-factor gap decomposition: why divergence happens}

Table~\ref{tab:gap-decomposition} places every factor's dominant
gap-decomposition term side by side, answering directly whether the same
configuration field drives disagreement across factors or whether the
dominant cause is factor-specific. This table is itself the primary
evidence for this paper's central claim -- that HXZ-versus-C\&Z divergence
has attributable, field-level causes -- rather than a supporting appendix
item.
% T4a/T4b: paper-anchored three-term decomposition + field-level gap
% decomposition, cross-factor. Source: comparison.json gap_decomposition +
% the CZ-P / HXZ-P three-term identity (Ch.3 eq. three-term).
\begin{table}[htbp]
  \centering
  \caption{Paper-anchored three-term decomposition, by factor.}
  \label{tab:three-term}
  \begin{tabular}{lcccc}
    \toprule
    Factor & Total gap ($CZ-P$) & Signal/env.\ term & Config term & Agent residual \\
    \midrule
    AssetGrowth & TODO & TODO & TODO & TODO \\
    GP          & TODO & TODO & TODO & TODO \\
    PS          & TODO & TODO & TODO & TODO \\
    BrandInvest & TODO & TODO & TODO & TODO \\
    OperProfRD  & TODO & TODO & TODO & TODO \\
    grcapx3y    & TODO & TODO & TODO & TODO \\
    \bottomrule
  \end{tabular}
\end{table}

\begin{table}[htbp]
  \centering
  \caption{Field-level gap decomposition, dominant term by factor.}
  \label{tab:gap-decomposition}
  \begin{tabular}{llccc}
    \toprule
    Factor & Dominant field & Contribution & Interaction terms & Explained fraction \\
    \midrule
    AssetGrowth & TODO & TODO & TODO & TODO \\
    GP          & Breakpoint source & TODO & TODO & TODO \\
    PS          & TODO & TODO & TODO & TODO \\
    BrandInvest & Weighting rule & TODO & TODO & TODO \\
    OperProfRD  & TODO & TODO & TODO & TODO \\
    grcapx3y    & TODO & TODO & TODO & TODO \\
    \bottomrule
  \end{tabular}
\end{table}

\begin{figure}[htbp]
  \centering
  \fbox{\parbox{0.9\linewidth}{\centering \vspace{4cm} F3: per-config-field contribution tornado plot \vspace{4cm}}}
  \caption{Field-level gap decomposition, all factors in the sample.}
  \label{fig:tornado}
\end{figure}

\subsection{$t$-channel decomposition}

Table~\ref{tab:t-channels} reports, for each factor and track pair, the
three log-$t$ channels of Equation~\eqref{eq:t-channel} -- how much of a
$t$-statistic change is attributable to the mean-return channel, the
volatility channel, or the sample-size channel.
% T5: log-t three-channel decomposition. Source: t_channel_decomposition.
\begin{table}[htbp]
  \centering
  \caption{$\log t$ decomposition into mean / volatility / sample-size channels.}
  \label{tab:t-channels}
  \begin{tabular}{lcccc}
    \toprule
    Factor & Track & Mean channel & Vol.\ channel & Sample-size channel \\
    \midrule
    TODO & TODO & TODO & TODO & TODO \\
    \bottomrule
  \end{tabular}
\end{table}


\subsection{Sensitivity band vs.\ HXZ and paper point estimates (headline result)}
\label{sec:sensitivity}

Table~\ref{tab:sensitivity-band} and Figure~\ref{fig:tband} together report
this paper's headline result: for each factor, the range of $t$-statistics
reachable across the menu of configurations $\Aref$ admits, plotted against
the paper's own reported $t$ and the $\HXZ$-standardized $t$. Where the
paper's point falls outside this band, or the $\HXZ$ point falls inside it,
the figure alone states the factor's conclusion without further argument.
% T6: sensitivity band (menu-reachable t min/max) vs. HXZ point vs. paper point.
% Source: robustness_summary + per-track t_stat. This is the flagship result
% (paired with figure F4).
\begin{table}[htbp]
  \centering
  \caption{Sensitivity band vs.\ HXZ and paper point estimates.}
  \label{tab:sensitivity-band}
  \begin{tabular}{lcccc}
    \toprule
    Factor & $t$ band (min) & $t$ band (max) & $t$ (HXZ) & $t$ (paper) \\
    \midrule
    AssetGrowth & TODO & TODO & TODO & TODO \\
    GP          & TODO & TODO & TODO & TODO \\
    PS          & TODO & TODO & TODO & TODO \\
    BrandInvest & TODO & TODO & TODO & TODO \\
    OperProfRD  & TODO & TODO & TODO & TODO \\
    grcapx3y    & TODO & TODO & TODO & TODO \\
    \bottomrule
  \end{tabular}
\end{table}

\begin{figure}[htbp]
  \centering
  \fbox{\parbox{0.9\linewidth}{\centering \vspace{5cm} F4: per-factor $t$-stat disagreement band vs.\ paper/HXZ points (flagship figure) \vspace{5cm}}}
  \caption{For each factor, the band spans the menu-reachable $t$-statistics;
    markers show the paper's own reported $t$ and the HXZ-standardized $t$.}
  \label{fig:tband}
\end{figure}

\subsection{Post-publication decay}
\label{sec:postpub}

The \texttt{insamp}/\texttt{between}/\texttt{postpub} split
(\S\ref{sec:pipeline}) is applied to test whether a factor's spread decays
after its publication date. [TODO: confirm per-factor sample coverage
before reporting this split -- fgr5yrLag's IBES-based sample and any factor
with a short overall window may not have enough post-publication months to
support the \texttt{between}/\texttt{postpub} comparison; report the exact
factor count used for this specific test explicitly if it differs from the
count used elsewhere in this section, rather than omitting the affected
factor silently.]

\subsection{Cross-factor outcome classification}

Table~\ref{tab:outcomes} closes the results section by classifying every
factor's replication outcome, completing the connection back to
\S\ref{sec:intro}'s contributions.
% T7: cross-factor replication outcome classification. Source: step7/step8
% classification labels (replicated_robustly / direction_only_replication /
% underidentified / etc.).
\begin{table}[htbp]
  \centering
  \caption{Cross-factor replication outcome classification.}
  \label{tab:outcomes}
  \begin{tabular}{ll}
    \toprule
    Factor & Outcome label \\
    \midrule
    AssetGrowth & TODO \\
    GP          & TODO \\
    PS          & TODO \\
    BrandInvest & TODO \\
    OperProfRD  & TODO \\
    grcapx3y    & TODO \\
    \bottomrule
  \end{tabular}
\end{table}


